\section{Termodinamikaning birinchi qonuni}



\begin{enumerate}

\item Termodinamika 1-qonuni nimani tavsiflaydi?
\begin{enumerate}[label=\Alph*)]
\item energiyaning saqlanish qonunini.
\item elastik deformatsiya energiyasini.
\item issiqlik muvozanatini.
\item issiqlik uzatilish jarayonini.
\end{enumerate}

\item Tashqi kuchlar gaz ustida 300 J ish bajarishdi. Agar bunda gazning ichki energiyasi 400 J ga ortgan bo‘lsa, unga qanday (J) issiqlik miqdori berilgan?
\begin{enumerate}[label=\Alph*)]
\item -100
\item 500
\item 100
\item 700
\end{enumerate}

\item Tashqi kuchlarga qarshi 250 J ish bajarildi va bunda gazning ichki energiyasi 550 J ortdi. Gazga qancha issiqlik miqdori berilgan (J)?
\begin{enumerate}[label=\Alph*)]
\item -300
\item 500
\item 800
\item 300
\end{enumerate}

\item Tashqi kuchlar gaz ustida 400 J ish bajardi va bunda gazning ichki energiyasi 550 J ortdi. Gazga qancha issiqlik miqdori berilgan (J)?
\begin{enumerate}[label=\Alph*)]
\item -800
\item 500
\item 800
\item 150
\end{enumerate}

\item Gazga 900 J issiqlik berish natijasida uning ichki energiyasi 725 J ortdi. Tashqi kuchlarga qarshi bajarilgan ishni toping (J).
\begin{enumerate}[label=\Alph*)]
\item -175
\item 500
\item 800
\item 175
\end{enumerate}

\item Gazga 900 J issiqlik berildi va gaz ustida 175 J ish bajarildi. Uning ichki energiyasi o‘zgarishini toping (J).
\begin{enumerate}[label=\Alph*)]
\item -725
\item 1075
\item -1075
\item 725
\end{enumerate}

\item Termodinamik jarayonda ideal gazning ichki energiyasi 300 J ga ortib, tashqariga 200 J issiqlik chiqqan. Bu jarayonda gaz qanday ish bajargan?
\begin{enumerate}[label=\Alph*)]
\item 500
\item -500
\item 100
\item -100
\end{enumerate}

\item Sistema ichki energiyalari teng, bir xil gaz bilan to‘ldirilgan va o‘zaro havosiz bo‘shliq bilan ajratilgan uchta idishdan iborat. Idishlardagi gazlarning massasi va harorati mos holda $m_1 = 4 \, kg, m_2 = 3 \, kg, m_3 = 0,5 \, kg$ va $T_1 = 300 \, K, T_2 = 400 \, K, T_3 = 600 \, K$. Vaqt o‘tishi bilan idishlardagi harorat qanday o‘zgaradi?
\begin{enumerate}[label=\Alph*)]
\item birinchi va ikkinchi idishlarda ko‘tariladi, uchinchisida pasayadi
\item ikkinchi va uchinchi idishlarda pasayadi, birinchisida ortadi
\item idishlardagi harorat o‘zgarmaydi
\item birinchi idishda ortadi, ikkinchida o‘zgarmaydi, uchinchisida pasayadi
\end{enumerate}

\item Qaysi qonun tabiatdagi jarayonlarning yo‘nalishini ko‘rsatadi?
\begin{enumerate}[label=\Alph*)]
\item energiyaning saqlanish qonuni
\item termodinamikaning I-qonuni
\item termodinamikaning II-qonuni
\item zaryad miqdorining saqlanish qonuni
\end{enumerate}

\item Kislorod izotermik kengayganda $A$ ish bajardi. Gaz bunda qancha issiqlik miqdori olgan?
\begin{enumerate}[label=\Alph*)]
\item $Q = A + \Delta U$
\item $Q = A$
\item $Q = -A$
\item $Q = A - \Delta U$
\end{enumerate}

\item Qanday jarayonda gazga berilgan issiqlik miqdorining hammasi bajarilgan ishga sarflanadi?
\begin{enumerate}[label=\Alph*)]
\item adiabatik
\item izobarik
\item izoxorik
\item izotermik
\end{enumerate}

\item Gaz izotermik kengayib 10 kJ ish bajardi. Gazning ichki energiyasi o‘zgarishini toping (kJ).
\begin{enumerate}[label=\Alph*)]
\item 0
\item 10
\item -10
\item 30
\end{enumerate}

\item Gaz izotermik kengayib 10 kJ ish bajardi. Gazga berilgan issiqlik miqdorini toping (kJ).
\begin{enumerate}[label=\Alph*)]
\item 0
\item 10
\item -10
\item 15
\end{enumerate}

\item Gaz izotermik kengayib 30 kJ ish bajardi. Gazga berilgan issiqlik miqdorini toping (kJ).
\begin{enumerate}[label=\Alph*)]
\item 0
\item 30
\item -10
\item 15
\end{enumerate}



\end{enumerate}







\begin{table}[!ht]
\centering
\begin{tabular}{|c|c|c|c|c|c|c|c|c|c|c|c|c|c|}
\hline
№ & J. & № & J. & № & J. & № & J. & № & J.& № & J. & № & J. \\ \hline
1 & A & 3 & C & 5 & D & 7 & B & 9 & C & 11 & D & 13 & B\\ \hline
2 & C & 4 & D & 6 & B & 8 & B & 10 & B & 12 & A & 14 & B\\ \hline
\end{tabular}
\end{table}
